\section{Q\&A：把 architecture slogan 还原成工程判断}

主讲结束后约 23 分钟 Q\&A 没有新增 slide，摄像机始终拍摄讲者与 conclusion 背景。它却包含本讲最重要的边界条件：specialization 如何保持 fusion、objective 与 parameters 如何分开审计、哪些想法只是值得实验，以及语言是否必然是 reasoning medium。本章把这些回答按工程问题重组，而不是逐问逐答抄写。

\subsection{参数分开后，信息怎样融合}

MoT 的 modality-specific projection 并没有取消 mixed-sequence attention。设 text query 为 $q_T=W_Q^T h_T$，image key 为 $k_V=W_K^V h_V$，仍可计算
$$
\operatorname{score}(T,V)=\frac{q_T^\top k_V}{\sqrt{d}}.
$$
$W_Q^T$ 与 $W_K^V$ 属于不同参数族，但 dot product 仍把两种 modality 放进同一 attention graph。模型因此可以“专用地编码，联合地读取”。

\teachervoice{00:42:42--00:44:08，老师用 self-attention 回答“参数隔离是否阻断信息交换”：token 经过不同参数，并不妨碍它们在 attention 中互相看见。读 architecture 图时应分别追踪 parameter path 与 information path。}

\begin{table}[H]
\centering
\small
\begin{tabularx}{0.98\textwidth}{L{0.23\textwidth}Y Y}
\toprule
设计轴 & 可选项 & 核心问题 \\
\midrule
Representation & VQ code、semantic embedding、VAE latent、audio codec & 输入/输出保留了什么信息？ \\
Objective & CE、diffusion、flow matching、action loss & 模型被奖励预测什么？ \\
Parameter sharing & fully shared、modality-specific、adapter、MoE & 哪些 token 竞争同一 capacity？ \\
Information mixing & global attention、cross-attention、late fusion & modality 在哪一层交换信息？ \\
Serving & synchronous、streaming、background reasoning & latency 与用户交互如何组织？ \\
\bottomrule
\end{tabularx}
\caption{多模态系统中必须分开审计的设计轴}
\end{table}

\subsection{Objective choice 与 parameter sharing 不是一回事}

一个 shared backbone 可以同时优化 CE 与 diffusion；一套 modality-specific 参数也可以输出 text。Q\&A 的价值在于拒绝“看到共享就假设同一 loss”“看到不同 loss 就假设两个模型”的二分法。工程实现应把每个 batch 的 modality mask、loss normalization、gradient norm 与 activated parameters 记录下来，才能定位 interference 来自哪里。

\begin{importantbox}{多任务训练的最小观测面}
每个 modality 分别记录 sample count、token count、loss、gradient norm、active parameters、吞吐与 validation metric；再记录跨 modality 的 attention/representation probe。没有这些分解，global training loss 下降可能掩盖某个 modality 持续退化。
\end{importantbox}

\subsection{Video generation 是否就是 world model}

Q\&A 讨论生成未来 frame 是否能帮助 action prediction。若模型能在条件动作下预测 $o_{t+1}$，它确实学到一部分 environment dynamics；但“看起来合理的视频”与“用于控制的 calibrated transition model”仍有差距。后者必须在 counterfactual action、rare failure 和 long-horizon rollout 上保持一致。

\teachervoice{00:46:16--00:48:00，老师把 video generation 视为可能的 predictive feedback，而不是直接宣布生成模型就是可靠 world model。讲义保留这一条件语气：需要任务反馈与闭环验证。}

\begin{warningbox}{World-model claim 的三项检查}
一看 action conditioning：不改变动作时预测是否相同？二看 rollout consistency：误差会否随时间爆炸？三看 decision utility：使用预测是否真的提高 policy success？若只有漂亮未来帧而没有这三项，最多证明视频生成能力。
\end{warningbox}

\subsection{“把文字画成图片”是实验问题，不是现成结论}

有同学问能否把 text 渲染成黑白 bitmap，从而用统一视觉 representation 建模一切。讲者明确说没有亲自做过，并判断 clean symbolic token 可能更高效。这个回答值得保留，因为它示范了对新颖想法的正确态度：先定义对照实验，而不是把直觉包装成事实。

一个合理实验应固定语料、model FLOPs 与 context information，比较：原始 text token；bitmap patch；OCR-pretrained bitmap encoder；混合 token。评价除 perplexity 外，还要看序列长度、训练吞吐、reasoning transfer 和 robustness。Bitmap 可能增加视觉共享，却也引入字体、渲染和空间冗余。

\teachervoice{00:48:14--00:53:31，老师反复标记“never tried it myself”“worth trying experiments”，最后只给出 efficiency intuition。讲义不把这个回答升级为 scaling law。}

\subsection{Object-centric representation 与 hierarchy}

Patch grid 方便计算，却不等于世界由固定方块组成。Q\&A 提出 object-oriented embedding、层级 representation 与更语义化的 scene unit。它们试图让模型表示“物体—关系—事件”，减少对像素级变化的敏感性。开放难题是如何在无监督数据中发现稳定对象、处理遮挡与身份连续性，并同时支持 generation。

\begin{knowledgebox}{Representation frontier 的三层结构}
低层保留纹理、声音频率和运动；中层形成 object、part、speaker、event；高层形成关系、意图和 task state。理解任务偏好高层 invariance，生成任务又需要低层 detail。一个可扩展系统很可能使用 hierarchy，而不是在单一 resolution 上解决所有任务。
\end{knowledgebox}

\subsection{Language 是当前 reasoning skeleton，但不是逻辑必需品}

最后一个长问题指出：文本 synthetic data 易生成、可隐藏 chain-of-thought，而视频 reasoning frames 昂贵且用户会直接看到。讲者同意 language 目前是非常有效的 abstraction skeleton，也同意视觉生成受效率和 UX 限制；但她没有排除未来纯视觉空间中的 reasoning。若模型准确预测复杂视频未来并支持决策，那也可视为一种非文本推理。

\teachervoice{01:03:32--01:04:21，老师的最终立场是“双重限定”：language skeleton 现在很有帮助；pure visual reasoning 在原则上仍是 open question。不要把当前工程优势改写成认知理论定律。}

\begin{warningbox}{空间推理回答的专业边界}
00:58:34，老师明确说自己不是 spatial reasoning 专家，只概括 robotics 近期进展。讲义因此不把这段回答作为 state-of-the-art survey；它只支持课程级判断：数字多模态能力与可靠空间推理之间仍有显著距离。
\end{warningbox}

\subsection{本章小结}

Q\&A 把“统一”和“原生”拆回工程变量：参数可以专用而信息仍融合；loss 与 backbone 是独立轴；video generation 只有在决策中有效才接近 world model；新 representation 假说需要 matched experiment；language 是当前高效的 reasoning interface，但未来形式仍开放。

